\section{Introduction}

A complete family of pairwise annihilating idempotents decomposes a module into exact coordinate pieces.  Many problems require more than one such decomposition.  A signal may be split first into two residue classes and then into four.  A filtered operator may be resolved into transfers that appear at different valuation depths.  A rational spectral splitting may exist over the fraction field while remaining invisible on the original lattice.  These are all refinement problems, but they involve different notions of scale.

The first part of the article concerns exact refinement.  If a fine packet family refines a coarse family, then each coarse idempotent is the sum of its fine descendants.  The coordinate spaces at the two levels are canonically isomorphic because both represent the whole module.  The refinement relation adds more information.  It tells us which fine coordinates belong to each coarse coordinate, and it gives a compatible decomposition of every operator block.  This is the algebra behind polyphase coordinates and nested filter banks.  Standard treatments of those subjects may be found in \cite{Vaidyanathan1993,Daubechies1992}.

The second part studies inversion in a complete filtered algebra.  Suppose an operator has the form
\[
T=D+N,
\]
where $D$ is invertible and every packet transfer in $N$ raises the filtration by a positive amount.  The ordinary Neumann series can then be expanded into weighted paths.  The weight is not the number of edges.  It is the total filtration gain along the path.  At precision $L$, every path of weight at least $L$ vanishes.  A finite transfer graph also gives a rational generating function for the number of paths of each total weight.  The relevant generating function facts are classical and are discussed in \cite{Stanley2012}.

The third part concerns integral separation.  Let $V$ be a vector space over the fraction field of a discrete valuation ring $\OO$, and let $\Lambda$ be an $\OO$ lattice in $V$.  A rational projector $E$ need not preserve $\Lambda$.  Its separation weight is the least $s$ such that
\[
\pi^sE\Lambda\subseteq\Lambda.
\]
A Bézout construction gives an upper bound through a polynomial resultant.  That bound need not be exact because it does not see the position of the lattice.  The exact answer is controlled by the finite module
\[
G_{\Lambda,\Pcal}
=
\Lambda\Big/\bigoplus_{\alpha}(\Lambda\cap V_\alpha),
\]
where $V=\bigoplus_\alpha V_\alpha$ is the rational packet decomposition.  The largest Smith exponent of this gluing module is exactly the common denominator needed to make every packet projector integral.  This interpretation also proves that separation weight cannot decrease under refinement.  Background on local fields, Smith normal form, and matrix polynomials may be found in \cite{Serre1979,Newman1972,GohbergLancasterRodman2009}.

The last part records how a defect between fine and coarse operators propagates.  If
\[
\Delta=CT_f-T_cC,
\]
then exact formulas describe the corresponding defect for powers, resolvents, holomorphic functions, exponentials, and a chain of several refinement levels.  These formulas are elementary, but collecting them in one place makes the later analytic uses transparent.  The spectral and resolvent background follows the usual operator theory in \cite{Kato1995}.

Every theorem in this paper concerns linear operators. An exact refinement retains the full object and changes its coordinates. Questions about nonlinear operation trees, nonlinear correction, and obstruction complexes require a separate theory.

\section{Four kinds of depth}\label{sec:four_depths}

Several quantities below are measured by nonnegative integers.  They do not measure the same thing.  Keeping them separate makes the later formulas easier to read.

The first is packet depth.  A finer packet family has more idempotents and a more detailed address set.  The passage from residues modulo two to residues modulo four is an example.  Packet depth is combinatorial.  It says how many distinctions have been made in the coordinate system.

The second is filtration depth.  A transfer in $F^r\Bcal$ is invisible modulo $F^r\Bcal$ and becomes visible at greater precision.  Two paths with the same number of edges can have different filtration depths.  A longer path can appear earlier when its edges have smaller weights.

The third is separation depth.  A rational packet projector may require multiplication by $\pi^s$ before it preserves the chosen lattice.  This number depends on both the rational decomposition and the lattice.  It can be zero even when the corresponding eigenvalues are close, and it can be positive even when the rational projector is simple.

The fourth is compatibility depth.  A fine operator and a separately chosen coarse operator need not commute with coarsening.  Their discrepancy is
\[
\Delta=CT_f-T_cC.
\]
This is not a packet depth or a denominator.  It measures the failure of one operator model to represent another exactly.

The four notions meet in applications, but no one of them determines the others.  An exact refinement tower may have zero coordinate error while its fine spectral packets have positive lattice weight.  A filtered inverse may truncate exactly at a given depth while the coarse operator still has a nonzero compatibility defect.  The paper treats each quantity with its own theorem before combining them.

\begin{keyidea}
Packet depth changes the address of a component. Filtration depth controls when a transfer survives a precision cutoff. Separation depth clears a rational denominator on a chosen lattice. Compatibility depth measures a discrepancy between two operator models. These quantities can interact, but each needs its own hypothesis and its own calculation.
\end{keyidea}

A simple two packet model already contains all four ideas.  Let
\[
X=X_0\oplus X_1,
\qquad
D=\begin{pmatrix}D_0&0\\0&D_1\end{pmatrix},
\qquad
N=\begin{pmatrix}0&N_{01}\\N_{10}&0\end{pmatrix}.
\]
A finer family may split $X_0$ and $X_1$ into smaller summands.  The blocks $N_{01}$ and $N_{10}$ may enter at different filtration weights.  Over a fraction field, the rational packet spaces may fail to split the chosen lattice.  A coarse operator on one selected summand may also fail to intertwine with the full matrix.  The later sections isolate these four effects one at a time.

